\scanpage{39}
\uline{Supposons que les car. résiduelles de $A$} \uline{soient \unsure{égales|positives}:} $p>0$
\[
\deleted{(x^{\unsure{(n)}})^{(m)}=}\qquad\text{op. de puiss. div. }(\pi_n)_{n\geq0}
\]
a)
\[
\pi_m^{\deleted{(m)}}(a)\pi_n^{\deleted{(n)}}(a)=\binom{m+n}{n}\pi_{m+n}^{\deleted{(m+n)}}(a)
\]
Supposons que $\left\{\begin{aligned}m&\equiv0\ (p)\\0&\leq n\leq p-1\end{aligned}\right.$ alors
\[
\deleted{\unsure{\pi^m}}\quad\binom{m+n}{n}=\frac{(m+1)\cdots(m+n)}{n!}\not\equiv0\ (p)
\]
donc
\[
\boxed{\pi_{m+n}^{\deleted{(m+n)}}(a)=\binom{m+n}{n}^{-1}\pi_m^{\deleted{(m)}}(a)\pi_n^{\deleted{(n)}}(a)}
\]
b)
\[
\pi_m^{\deleted{m}}\circ\pi_n^{\deleted{n}}=\frac{(mn)!}{(n!)^m m!}\pi_{mn}^{\deleted{mn}}
\]
[ Supposons $m=p$, on trouve un coefficient
\[
c_n=\frac{(pn)!}{(n!)^p p!}
\]
et on a, si $n=a_0+a_1p+\cdots+a_rp^r$
\[
\begin{aligned}
\nu_p(c_n)&=\deleted{\unsure{a}}\ \nu_p((pn)!)-p\nu_p(\unsure{m|n}!)-1\\
&=\frac{pn-\sum a_i}{p-1}-p\frac{n-\sum a_i}{p-1}-1\\
&=\deleted{\frac{1}{p-1}}(\sum a_i)-1
\end{aligned}
\]
]
\figure{\begin{tikzpicture}\draw[line width=1.1pt] (0,0)--(0.18,0.3);\end{tikzpicture}}
Supposons $n=p$, on trouve un coeff.
\[
d_m=\frac{(pm)!}{(p!)^m m!}
\]
\margin{$m=\deleted{\sum}b_0+b_1p+\cdots+b_sp^s$}
et
\[
\begin{aligned}
\nu_p(d_m)&=\nu_p((pm)!)-m\nu_p(p!)-\nu_p(m!)\\
&=\frac{pm-\sum b_i}{p-1}-m-\frac{m-\sum b_i}{p-1}\\
&=0
\end{aligned}
\]
donc
\[
\pi_{pm}^{\deleted{pm}}=\left(\frac{(pm)!}{(p!)^m m!}\right)^{-1}\pi_m^{\deleted{m}}\circ\pi_p^{\deleted{p}}
\]
a) et b) ensemble nous montrent que \uline{les} $\pi_n^{\deleted{n}}$ \uline{sont connus} \uline{quand on connaît} $\pi_p^{\deleted{p}}$
\mspage{38}
